% 3-4-features.tex
%  Budget: 7pp -- the largest section of chapter 3.
%  Source map: the LEAR feature convention (price lags D-1/D-2/D-3/D-7, exog
%   lags D-1/D-7, same-day exog forecasts, weekday dummies) and the fairness
%   argument for reusing it -> lago_forecasting_2021 ;
%   earlier DNN/LEAR feature treatment in the same benchmark line
%   -> lago_forecasting_2018 ;
%   fuel prices as a documented driver this thesis deliberately excludes
%   -> chai_forecasting_2024 , trebbien_explainable_2023 .
%  Matrix verified: (2177, 247) = 96 price-lag + 96 exog-lag + 48 exog-D0 + 7 dow.
%  CRITICAL: 3-4 must justify why exog_*_D0 is legal (day-ahead forecasts
%   published BEFORE the origin, not realized values). HANDOFF section 10.
%  NUMERALS: integers -> Persian digits; decimals/thousands -> \lr{} (Latin).
%   Verified on render 2026-09-05; Persian decimal marks are absent from
%   B Nazanin and reorder. See fix_numerals.py.

\section{مهندسی ویژگی}
\label{sec:3-4}

\subsection*{تنظیم مسئله}

صورت‌بندی مسئله در این پژوهش چنین است: در مبدأ پیش‌بینی، یعنی روز \lr{O} و پیش از
بسته‌شدن حراج روز-پیش، هر بیست‌وچهار قیمت ساعتی روز \lr{O+1} پیش‌بینی می‌شود. خروجی
هر مدل یک بردار ۲۴مؤلفه‌ای است، نه یک عدد؛ همین انتخاب، افق پیش‌بینی را با ساختار
واقعی حراج هم‌شکل می‌کند، چون در بازار روز-پیش نیز هر بیست‌وچهار ساعت به‌صورت
یک‌جا و در یک حراج تسویه می‌شوند.

ماتریس ویژگیِ حاصل ۲۱۷۷ سطر و ۲۴۷ ستون دارد. هر سطر یک روز هدف است و هر ستون یک
ویژگی. این ماتریس، منبعِ یکتای ویژگی برای هر پنج مدل پژوهش است: مدل ساده، \lr{SARIMAX}،
\lr{LEAR-LASSO}، \lr{LightGBM} و \lr{LSTM} همگی از همین ماتریس تغذیه می‌شوند و هیچ‌یک
مسیر ویژگی‌سازیِ اختصاصی ندارد.

یکتاییِ ماتریس نه از سرِ سهولتِ مهندسی، که تصمیمی روش‌شناختی است. اگر هر مدل ماتریس خودش را
می‌ساخت، دو هزینه پدید می‌آمد. نخست، هر مسیر تازه سطح تازه‌ای برای نشت اطلاعات
می‌گشود که آزمون‌های این بخش آن را پوشش نمی‌دادند. دوم، و مهم‌تر، بخشی از اختلاف
دقتِ فصل چهارم دیگر به خودِ مدل نسبت‌دادنی نبود؛ می‌توانست از تفاوت در ورودی
برخاسته باشد. مقایسه‌ی منصفانه ایجاب می‌کند که تنها یک متغیر میان مدل‌ها تغییر کند، و
آن متغیر خودِ مدل باشد.

\subsection*{قید علّی: خط قرمز پژوهش}

قاعده‌ی حاکم بر کل این بخش یک جمله است: \emph{هیچ ویژگی‌ای نباید از اطلاعات پس از
مبدأ پیش‌بینی استفاده کند.}

نقض این قاعده، شایع‌ترین و مخرب‌ترین خطای این حوزه است، چون بی‌سروصدا رخ می‌دهد. مدلی
که یک ستون آلوده دریافت کرده باشد، بدون هیچ خطا یا هشداری آموزش می‌بیند، اعداد دقتِ
چشمگیری تولید می‌کند، و آن اعداد در عمل بی‌معنا هستند. تفاوت میان یک پژوهش معتبر و
یک پژوهش بی‌اعتبار در این حوزه، اغلب یک ستون است.

به همین دلیل این قاعده در پروژه به‌صورت قرارداد رعایت نمی‌شود، بلکه با دو سازوکار
اجرایی اعمال می‌گردد.

\textbf{سازوکار نخست: اعلان تصریحیِ در دسترس بودن.} هر ستون ماتریس ویژگی باید در یک
سامانه‌ی ثبت، منبع خود و میزان وقفه‌اش نسبت به روز هدف را اعلام کند. سپس تابعی تصریحی
کل مجموعه‌ی ستون‌ها را می‌پیماید و در صورت یافتن ستونی که اعلان نشده باشد، یا ستونی که
اعلانش با نامش نخواند، اجرا را با خطا متوقف می‌کند. نکته‌ی طراحیِ مهم آنکه حالت
پیش‌فرض این وارسی سخت‌گیرانه است: ستونی که فراموش شده اعلام شود، مجاز فرض نمی‌شود،
بلکه مردود است. سامانه‌ای که ستون ناشناخته را در سکوت بپذیرد، دقیقاً همان چیزی است
که باید از آن محافظت کند.

\textbf{سازوکار دوم: آزمون خودکار.} یک آزمون در مجموعه‌ی آزمون‌های پروژه، این ادعا را
بر ماتریس واقعی و با پیکربندی واقعی می‌سنجد، نه بر نمونه‌ای ساختگی. آزمون دومی نیز
دسترسیِ کد مهندسی ویژگی به شبکه را ممنوع می‌کند. علت آن، هم بازتولیدپذیری است و هم
همین قید علّی: کدی که بتواند در زمان ساخت ویژگی به یک واسط زنده وصل شود، می‌تواند
مقداری متأخرتر از مبدأ پیش‌بینی را دریافت کند، بی‌آنکه در متن برنامه چیزی به آن
اشاره کند.

\subsection*{چهار خانواده ویژگی}

مجموعه‌ی ویژگی‌ها عیناً از قرارداد مدل مرجع \lr{LEAR} پیروی می‌کند
\cite{lago_forecasting_2021, lago_forecasting_2018}. این پیروی، انتخابی آگاهانه و
محدودکننده است: می‌شد مجموعه‌ی غنی‌تری ساخت، اما آنگاه مقایسه‌ی فصل چهارم با اعداد
منتشرشده دیگر هم‌مقیاس نبود و هر تفاوتی در دقت را می‌شد به تفاوت در ورودی نسبت داد،
نه به مدل. ساختار ۲۴۷ ستون به شرح زیر است.

\subsubsection*{وقفه‌های قیمت (۹۶ ستون)}

بردار کامل ۲۴ساعته‌ی قیمت برای چهار روز گذشته نسبت به روز هدف: \lr{D-1}، \lr{D-2}،
\lr{D-3} و \lr{D-7}. چهار روز ضربدر ۲۴ ساعت، ۹۶ ستون می‌شود.

انتخاب همین چهار وقفه \lr{—} و نه مثلاً پنج یا ده \lr{—} دو توجیه دارد که یکدیگر را
تقویت می‌کنند. توجیه نخست، هم‌ارزی با قرارداد مرجع است. توجیه دوم تجربی است و در \cref{sec:3-3-3} نشان داده شد: تابع خودهمبستگی قلّه‌های آشکاری در وقفه‌های ۲۴ و ۱۶۸ ساعت دارد
(به‌ترتیب \lr{0.6598} و \lr{0.5855})، یعنی روز پیشین و همان روز در هفته‌ی پیشین
بیشترین اطلاع را درباره‌ی روز هدف دارند. وقفه‌های \lr{D-2} و \lr{D-3} فاصله‌ی میان این
دو را پر می‌کنند و برای دوشنبه‌ها نقش ویژه دارند، چون آخرین روز کاریِ پیش از دوشنبه،
جمعه است.

\subsubsection*{وقفه‌های متغیرهای برون‌زا (۹۶ ستون)}

بردار ۲۴ساعته‌ی دو متغیر برون‌زا \lr{—} پیش‌بینی بار و پیش‌بینی تولید بادی و خورشیدی
\lr{—} برای دو وقفه‌ی \lr{D-1} و \lr{D-7}. دو متغیر ضربدر دو وقفه ضربدر ۲۴ ساعت، ۹۶
ستون.

\subsubsection*{متغیرهای برون‌زای خودِ روز هدف (۴۸ ستون)}

دو متغیر برون‌زا برای \emph{خودِ روز هدف}، یعنی روزی که قیمتش هنوز پیش‌بینی نشده
است. دو متغیر ضربدر ۲۴ ساعت، ۴۸ ستون.

این خانواده تنها جایی است که ماتریس ویژگی مقداری از روز هدف را در بر می‌گیرد، و
بنابراین تنها جایی است که خواننده حق دارد از نشت اطلاعات بپرسد. پاسخ در ماهیت این دو
سری است و باید صریح گفته شود.

\lr{exog\_1} و \lr{exog\_2} مقدارِ \emph{محقق‌شده‌ی} بار و تولید تجدیدپذیر نیستند؛
\emph{پیش‌بینی روزپیشِ} آن‌ها هستند، که بهره‌بردار شبکه پیش از بسته‌شدن حراج روز-پیش
منتشر می‌کند. یعنی در مبدأ پیش‌بینی، این اعداد پیشاپیش منتشر شده و در دسترس هر
شرکت‌کننده‌ی بازار هستند. استفاده از آن‌ها نه‌تنها نشت نیست، بلکه بازتاب همان
اطلاعاتی است که یک معامله‌گر واقعی در همان لحظه در اختیار دارد. مدلی که آن‌ها را
نادیده بگیرد، از اطلاعات مجازِ در دسترس استفاده نکرده است.

تمایز کلیدی میان «مقدار محقق‌شده‌ی روز هدف» و «پیش‌بینیِ منتشرشده‌ی روز هدف» است. اولی
نشت است، دومی نیست. سامانه‌ی اعلانِ در دسترس بودن نیز همین تمایز را رمزگذاری می‌کند:
این ستون‌ها با وقفه‌ی صفر اما با منبعِ «پیش‌بینیِ منتشرشده» اعلام می‌شوند، و ستونی که
مقدار محقق‌شده با وقفه‌ی صفر اعلام کند، مردود می‌شود.

\subsubsection*{متغیرهای مجازی روز هفته (۷ ستون)}

هفت متغیر مجازیِ دودویی برای روز هفته‌ی روز هدف. \cref{sec:3-3-3} شکاف بزرگی میان روزهای
هفته نشان داد \lr{—} میانگین یکشنبه \lr{23.33} در برابر \lr{38.70} برای سه‌شنبه
\lr{—} که وجود این ستون‌ها را توجیه می‌کند. با این‌همه، تحلیل تفسیرپذیریِ \cref{sec:4-6}
نشان خواهد داد که سهم آن‌ها در حضور وقفه‌های قیمتی ناچیز است؛ یعنی اطلاعات روز هفته
پیشاپیش در خودِ قیمت‌های گذشته رمزگذاری شده است.

جمع چهار خانواده: ۹۶ به‌علاوه‌ی ۹۶ به‌علاوه‌ی ۴۸ به‌علاوه‌ی ۷ برابر ۲۴۷ ستون. این تطابق
عددی با فرمول ویژگیِ خودِ کتابخانه‌ی مرجع، در کد راستی‌آزمایی می‌شود و صرفاً ادعا
نمی‌شود.

\subsection*{آنچه عمداً در مجموعه ویژگی نیست}

چند متغیر که ادبیات آن‌ها را مؤثر می‌داند در این ماتریس نیستند، و تصریح دلیل آن،
بخشی از صداقت روش‌شناختی این پژوهش است.

\textbf{قیمت سوخت.} پژوهش‌ها نشان داده‌اند که قیمت حامل‌های انرژی از عوامل تعیین‌کننده‌ی
قیمت برق است و لحاظ‌کردن آن می‌تواند دقت را بهبود دهد
\cite{chai_forecasting_2024, trebbien_explainable_2023}. با این‌حال، سری‌های مرجعِ
گاز و زغال‌سنگ تجاری و دارای مجوزند و در چارچوب بازتولیدپذیریِ این پژوهش قابل
استفاده نبودند؛ منابعی هم که سری‌های \emph{مشتق‌شده} از آن‌ها را بازنشر می‌کنند، در
این پژوهش به‌عنوان مرجع استناد به کار می‌روند و نه به‌عنوان منبع داده. پیامد مستقیم
این محدودیت آن است که شاخص‌های حاشیه‌ی سوخت \lr{—} مانند اسپرد تمیز \lr{—} در این
پایان‌نامه ساختنی نبودند.

\textbf{هواشناسی و جریان‌های مرزی.} این دو نیز عمداً کنار گذاشته شده‌اند، اما به دلیلی
متفاوت: افزودن آن‌ها هم‌مقیاسیِ مقایسه با اعداد منتشرشده‌ی مرجع را از میان می‌برد. این
تصمیم، مبادله‌ای آگاهانه است و نه غفلت. پژوهش، دقتِ ممکن را در ازای امکانِ مقایسه‌ی
معتبر محدود کرده است، و \cref{sec:5-4} افزودن این متغیرها را به‌عنوان مسیر طبیعی پژوهش‌های
آینده پیشنهاد می‌کند.

ساختار کد این امکان را حفظ کرده است: بلوک‌های ویژگیِ فیزیکی به‌صورت افزودنی و
به‌طور پیش‌فرض خاموش پیاده‌سازی شده‌اند، به‌طوری‌که مسیر پیش‌فرضِ ساخت ویژگی دقیقاً
همان مجموعه‌ی ۲۴۷ ستونی است که نتایج تثبیت‌شده‌ی فصل چهارم از آن تولید شده‌اند. بلوکی که
به داده‌ی ثبت‌نام‌محور نیاز دارد نیز، به‌جای جایگزین‌کردن یک سریِ تقریبی، با خطا متوقف
می‌شود؛ یعنی نبودِ داده در سکوت به یک عدد ساختگی بدل نمی‌گردد.

\subsection*{هدف روزانه}

افزون بر هدف ساعتی، پژوهش هدف دومی نیز دارد: میانگین بار پایه‌ی روزانه، یعنی میانگین
۲۴ قیمت ساعتی هر روز. این هدف از دو مسیر متفاوت ساخته می‌شود، و همین دوگانگی است که
پرسش پژوهشی چهارم را پاسخ می‌دهد.

مسیر نخست، \emph{مدل‌سازی مستقیم} است: هدف روزانه به‌عنوان یک متغیر تک‌مؤلفه‌ای
تعریف و مدل مستقیماً بر آن برازش می‌شود. مسیر دوم، \emph{تجمیع} است: مدلِ ساعتی
پیش‌بینی خود را تولید می‌کند و سپس از ۲۴ مؤلفه‌ی آن میانگین گرفته می‌شود.

برای آنکه مقایسه‌ی این دو مسیر معنادار بماند، هر دو از همان ماتریس ویژگی، همان مبادی
پیش‌بینی و همان بودجه‌ی تنظیم ابرپارامتر استفاده می‌کنند. اگر مسیر مستقیم بودجه‌ی
تنظیم کمتری می‌گرفت، هر برتریِ مسیر تجمیعی می‌توانست صرفاً بازتاب همان نابرابری
باشد. برابری بودجه، همان چیزی است که مقایسه‌ی \cref{sec:4-4} را از یک مصنوعِ طراحی مصون
می‌دارد.
